\subsection{Fundamentos y conceptos relevantes para el proyecto}\label{sec:fundamentos}
\textbf{Modelado físico del incendio.} FDS es el software de simulación que representa el incendio mediante la interacción entre combustible, transferencia de calor, combustión y movimiento del aire. El trabajo de Mell y colaboradores constituye un antecedente de la aproximación física a incendios de pastizales~\cite{mell2007}, mientras que Valero, Jofre y Torres estudian el uso de configuraciones de distinta fidelidad para analizar propagación, sensibilidad e incertidumbre~\cite{valero2021}. En su formulación CFD, FDS utiliza una representación de las ecuaciones de movimiento para flujos de bajo número de Mach impulsados por calor, por lo cual resuelve un problema tridimensional y transitorio. Las capacidades de representación de vegetación asociadas históricamente a WFDS se utilizan dentro de FDS según la convención de la sección~\ref{sec:solucion}~\cite{fdsnist,wfdsusfs}. Esta formulación se distingue de los modelos de propagación de frente.

\textbf{Discretización y costo de cálculo.} La malla CFD divide el dominio físico en celdas donde se aproximan las variables del flujo y sus intercambios. Si un dominio rectangular se discretiza de manera uniforme, su número de celdas puede expresarse como $N=(L_x/\delta x)(L_y/\delta y)(L_z/\delta z)$, donde las longitudes y los tamaños de celda deben utilizar unidades compatibles. Por tanto, reducir a la mitad el tamaño de celda en las tres direcciones multiplica geométricamente por ocho el número de celdas del mismo dominio, aunque el costo temporal real también depende de la integración, los modelos y la ejecución. Este razonamiento explica por qué la fidelidad requiere decisiones explícitas de resolución, pero no proporciona una medición de tiempo o memoria para el sector del proyecto. El análisis de Valero y colaboradores aporta un antecedente para estudiar esa relación y sus efectos en cantidades de interés~\cite{valero2021}.

\textbf{Combustibles y representación territorial.} Una cobertura vegetal cartográfica identifica clases y distribución espacial, pero no define por sí sola la estructura tridimensional ni todas las propiedades físicas que requiere una entrada de simulación. La preparación debe distinguir composición y geometría, carga, humedad y las propiedades pertinentes al modelo elegido, conservando sus unidades y condiciones de aplicación. En el contexto de los Cerros Orientales, Verano Velásquez caracteriza información territorial para FARSITE~\cite{verano2013}, y Méndez Álvarez estudia la calibración de modelos de combustible altoandino~\cite{mendez2018}. Sus aportes sirven para examinar fuentes, parámetros y dificultades locales, aunque la transferencia a FDS exige comprobar correspondencias y restricciones. Una calibración realizada para otro modelo no se convierte automáticamente en una calibración de FDS, ni una cobertura histórica equivale a un inventario actual y exacto de especies y combustibles.

\textbf{Amenaza y distribución espacial de igniciones.} Morales Salguero integra información territorial mediante álgebra de mapas y simulación para estudiar amenaza de ignición y propagación en los Cerros Orientales~\cite{morales2017}. Las categorías resultantes permiten distinguir niveles de amenaza, pero no deben interpretarse como probabilidades estadísticas calibradas de que ocurra una ignición. Esta distinción es relevante para fundamentar la distribución de referencia establecida en la sección~\ref{sec:solucion}, cuya regla de espaciamiento debe ser definida por el proyecto. Los puntos de ignición representan posiciones donde se preparan escenarios, mientras que las celdas CFD representan la discretización de cada experimento físico, de modo que aumentar puntos disponibles no equivale a refinar la resolución de las ecuaciones.

\textbf{Cómputo de alto desempeño y experimentos precalculados.} Existen dos niveles de paralelización que requieren comprobaciones diferentes: ejecutar escenarios independientes de forma concurrente y distribuir el cálculo de un escenario entre recursos que intercambian información. El primer nivel aprovecha la independencia de combinaciones de entrada, mientras que el segundo depende de la descomposición y comunicación admitidas por la revisión del simulador. La distribución de escenarios independientes y la distribución interna de un caso plantean, por tanto, necesidades diferentes de memoria, almacenamiento y comunicación. Esta distinción fundamenta las comprobaciones de viabilidad y las mediciones del piloto descritas en las secciones~\ref{sec:inicio-actividades} y~\ref{sec:construccion-actividades}.

\textbf{Consulta por semejanza.} Recuperar escenarios requiere representar la situación observada y los casos disponibles mediante variables comparables, con una definición explícita de cercanía y tratamiento de unidades. Cover y Hart ofrecen un antecedente fundacional sobre vecinos cercanos en clasificación~\cite{cover1967}, del cual puede retomarse el principio de utilizar proximidad entre representaciones para recuperar información. Sin embargo, sus resultados de clasificación no demuestran la precisión física de una búsqueda de incendios ni seleccionan automáticamente la métrica adecuada para este problema. La normalización de variables, la dirección del viento, la proximidad espacial y las condiciones no cubiertas deberán tratarse según las propiedades de la biblioteca, y la cercanía de un resultado deberá comunicarse como correspondencia entre condiciones, sin convertirla en garantía de representatividad.

\textbf{Incertidumbre, multifidelidad y agregación.} La incertidumbre puede proceder de datos de entrada, simplificaciones del modelo, discretización y diferencias entre las condiciones observadas y las disponibles. Monte Carlo estudia variabilidad mediante repeticiones con entradas muestreadas, mientras que los enfoques multinivel combinan estimaciones de distinta fidelidad para reducir el costo de estudiar ciertas cantidades estadísticas. El artículo seminal de Valero y colaboradores utiliza ese marco para analizar simulación de incendios~\cite{valero2021}. Su objeto es diferente del procedimiento espacial de la sección~\ref{sec:solucion}: la estimación estadística entre fidelidades no constituye una regla de combinación de frentes. La superposición de frentes procedentes de igniciones calculadas de forma independiente no reproduce necesariamente todas las interacciones de un incendio iniciado desde un perímetro, por lo que requiere un contraste propio. En consecuencia, multifidelidad y Monte Carlo multinivel se conservan como fundamentos bibliográficos y no se incorporan por esta referencia como componentes comprometidos del producto.

\textbf{Representación geográfica e interpretación.} Una visualización georreferenciada vincula los resultados a un sistema espacial y a una secuencia temporal coherentes con la información consultada. Su corrección exige comprobar coordenadas, unidades, orientación y correspondencia de tiempos, además de diferenciar lo observado de lo estimado. Estas propiedades permiten examinar que la animación represente los datos procesados, aunque no bastan para acreditar que el usuario comprenda su alcance. Por ello, la evaluación distinguirá la coherencia técnica de la representación y la comprensión de las tareas de uso, conforme a las fases de las secciones~\ref{sec:construccion} y~\ref{sec:transicion}.
